\subsection{Reduction to a Scalar Optimisation}
\label{sec:reduction}

The equal-spacing constraint fixes all $N$ target positions given a
single free variable, the \emph{reference phase} $\uref \in [0, 2\pi)$.
After sorting the satellites by their current argument of latitude
$u_1 \le u_2 \le \cdots \le u_N$, the $i$-th satellite's target is:
\begin{equation}
  u_{i,\text{target}}(\uref)
    = \uref + (i-1)\,\frac{2\pi}{N},
  \quad i = 1, \dots, N.
  \label{eq:target}
\end{equation}
The required angular transfer for each satellite is:
\begin{equation}
  \du_i(\uref)
    = \operatorname{wrap}_{[-\pi,\pi]}\!\bigl(u_{i,\text{target}}(\uref) - u_i\bigr),
  \label{eq:delta}
\end{equation}
where $\operatorname{wrap}_{[-\pi,\pi]}(x) = ((x + \pi) \bmod 2\pi) - \pi$
selects the shortest-path transfer for each satellite (at most half a revolution).

Substituting~\eqref{eq:delta} into the $\dv$ formula~\eqref{eq:dv} gives
a cost that depends on $\uref$ only. The full $N$-satellite scheduling problem
therefore reduces to a \textbf{scalar search over $\uref$}, solved
with \MATLAB's \texttt{fminbnd} on $[0, 2\pi)$.

\subsection{Objectives}
\label{sec:objectives}

Let $\dv_i(\uref)$ be the cost for satellite~$i$ from~\eqref{eq:dv},
and let $r_i(\uref) = f_i - \dv_i(\uref)$ be the remaining fuel. Two
objectives are defined:
\begin{align}
  J_1(\uref) &= \sum_{i=1}^{N} \dv_i(\uref)
    && \text{(total $\dv$ consumed by the fleet),}
    \label{eq:J1}\\
  J_2(\uref) &= \max_{i}\,\bigl|r_i(\uref) - \bar{r}(\uref)\bigr|
    && \text{(max deviation of remaining budgets from their mean).}
    \label{eq:J2}
\end{align}

\subsection{Optimisation Strategies}

\subsubsection{Weighted Strategy}

A single \texttt{fminbnd} call minimises the scalarised objective:
\begin{equation}
  \min_{\uref \in [0,2\pi)}\; W_1\,J_1(\uref) + W_2\,J_2(\uref).
  \label{eq:weighted}
\end{equation}
The weights $W_1$ and $W_2$ control the trade-off between economy and
balance. In the test case we use $W_1 = 1$ and $W_2 = 10$, placing
a strong penalty on fuel imbalance.

\subsubsection{Lexicographic Strategy}

Total fuel economy is treated as strictly more important than balance.
Two sequential \texttt{fminbnd} calls are used:
\begin{enumerate}
  \item \textbf{Stage~1:} minimise $J_1$ to obtain the optimal reference
        phase $\uref^*$ and the minimum total cost $J_1^* = J_1(\uref^*)$.
  \item \textbf{Stage~2:} minimise $J_2$ subject to
        $J_1(\uref) \le J_1^*(1 + \varepsilon)$, enforced via a large
        penalty term with $\varepsilon = 10^{-6}$.
\end{enumerate}
This guarantees that fuel balance is only improved when doing so does
not meaningfully increase total consumption.
